\section{The arithmetic Floquet normal form and critical recurrence law}\label{sec:floquet}
Put \[
 \alpha=1+\sqrt2,\quad \lambda=\alpha^2,\quad
 T=\begin{pmatrix}3&2\\4&3\end{pmatrix}.
\] Let K tend to infinity through odd primes, and define positive integers N,F by \[
 N+F\sqrt2=\alpha^K,\qquad N^2-2F^2=-1.
\] The moduli N may be composite. On the odd subspace of the standard N-dimensional torus quantization, let \(U_{N}\) be the phase-normalized quantization of T described in Subsection~\ref{sec:floquet}.1, and let \[
 B_N f(x)=2\cos(2\pi x/N)f(x),\quad
 U_{\varepsilon,N}=e^{i\varepsilon B_N}U_{N,\mathrm{Fl}},\quad
 d_N=(N-1)/2.
\]

\begin{theorem}[critical resonance law]\label{pass1:ch07_floquet_resonance-1} For every fixed real \(\tau\) and every integer m, with \(\varepsilon=\tau/\sqrt{2K}\), \[
 \lim_{K\longrightarrow\infty}
 \frac{1}{d_N}\operatorname{Tr}_{-}
       U_{\varepsilon,N}^{mK}
       =\exp(-m^2\tau^2/2).                                      \tag{FN.A}
\] Consequently the empirical eigenvalue measure of \(U_{\varepsilon,N}^K\) converges weakly on the unit circle to the wrapped normal law. For \(\tau\) nonzero its angular density on {[}-\(\pi\),\(\pi\)) is \[
 w_\tau(\theta)=\frac{1}{\sqrt{2\pi}|\tau|}
       \sum_{\ell\in\mathbb Z}
       \exp\!\left[-\frac{(\theta+2\pi\ell)^2}{2\tau^2}\right].   \tag{FN.B}
\] For \(\tau=0\) the limit is the point mass at 1. In particular the normalized resonance echo obeys \[
 d_N^{-2}\left|\operatorname{Tr}_{-}U_{\varepsilon,N}^{mK}\right|^2
       \longrightarrow e^{-m^2\tau^2}.                          \tag{FN.C}
\] This is a deterministic limit: there is no parameter average in (FN.A).
\end{theorem}

The principal finite estimate, proved before any limit is taken, is \[
 \left\|V_\varepsilon U_{\varepsilon,N}V_\varepsilon^*
            -e^{i\varepsilon\overline B}U_{N,\mathrm{Fl}}\right\|_{2,-}
       \leq \frac{12\varepsilon^2}{\sqrt K},\qquad K\geq11,        \tag{FN.D}
\] for an explicit unitary \(V_{\varepsilon}\). Here the norm is normalized Hilbert--Schmidt norm on the entire odd subspace. It does not control an individual level or an operator-norm spectral gap. The finite comparison error in (FN.A) is at most \(6|m|\tau^2/\sqrt K\); an additional Gaussian characteristic-function convergence error is \(o(1)\), not assigned that rate.

\subsection{Arithmetic and the fixed classical observable}\label{floquet:arithmetic-and-the-fixed-classical-observable}

The recurrence \((N_j,F_j)=(N_{j-1}+2F_{j-1},N_{j-1}+F_{j-1})\), starting at \((1,0)\), shows that \(N\) is odd. In \(\mathbb Z/N\mathbb Z\), put \[
 s=F^{-1},\qquad s^2=2,\qquad h=3+2s,\qquad q=-F.
\] Since \((1+s)^K=N+Fs=1\), we have \(h^K=1\). Also \(h-1=2(1+s)\) is a unit. Modulo every prime \(\ell\mid N\), the order of \(h\) is precisely \(K\): it divides the prime \(K\) and is not one. Hence \(1-h^j\) is a unit for \(1\leq j<K\), and \(1+h^j\) is a unit for every \(j\), because \(K\) is odd. In particular, \(\ell\equiv1\pmod K\) and \(N\equiv1\pmod{2K}\). Neither squarefreeness nor primality of \(N\) is required.

Use the counting-measure inner product and centered Weyl operators \[
 W(a,b)f(x)=e_N(a(x-b/2))f(x-b),\qquad
 S_u f(x)=f(u^{-1}x),\qquad Rf(x)=f(-x).
\] Divisions in Weyl exponents mean inverses in \(\mathbb Z/N\mathbb Z\). The multiplication and covariance formulas are \[
 W(a,b)W(c,d)=e_N((ad-bc)/2)W(a+c,b+d),
\] \[
 S_uW(a,b)S_u^*=W(u^{-1}a,ub),\qquad RW(v)R=W(-v).
\] Individual Weyl translations need not preserve parity; their symmetrized sums do. On the odd subspace take \[
 U_0=S_h^{-1},\qquad B=W(q,1)+W(q,1)^*,\qquad
 U_\varepsilon=e^{i\varepsilon B}U_0.                            \tag{FN.1}
\] In these coordinates \(U_0^K=I\) exactly.

For completeness, this is unitarily equivalent to the fixed-map formulation of Theorem~\ref{pass1:ch07_floquet_resonance-1}, including the kick, not merely its compressed observable. Let \(C_tf(x)=e_N(tx^2/2)f(x)\), and let \(F_N\) have kernel \(e_N(xy)/\sqrt N\). Set \[
 M=\begin{pmatrix}-1/(2s)&1\\-1/2&-s\end{pmatrix},\qquad
 V=C_{s^{-1}}S_{1/2}F_N C_{2s}.
\] Its factors have Weyl covariance matrices respectively \[
 \begin{pmatrix}1&t\\0&1\end{pmatrix},\quad
 \begin{pmatrix}u^{-1}&0\\0&u\end{pmatrix},\quad
 \begin{pmatrix}0&1\\-1&0\end{pmatrix}.
\] Multiplication gives covariance \(M\), \(\det M=1\), and \(T=M\operatorname{diag}(h,h^{-1})M^{-1}\). Define \(U_{N,\mathrm{Fl}}=VU_0V^*\). The index of \(V^*W(1,0)V\) is \(M^{-1}(1,0)=(-s,1/2)\). Further conjugation by \(S_2\), which commutes with \(U_0\), changes this to \((-s/2,1)=(-F,1)\pmod N\). Thus \(S_2V^*\) conjugates the entire physical kicked propagator to (FN.1). These unitaries all preserve parity.

With the usual semiclassical convention \(e^{i\,\mathrm{action}/\hbar}\), the classical kick amplitude associated with \(\varepsilon\) is proportional to \(\hbar\varepsilon\). Here \(\hbar\asymp1/N\). At the critical scale the amplitude is of order \(1/(N\sqrt K)\), not a fixed classical perturbation.

\subsection{Two summable phase estimates}\label{floquet:two-summable-phase-estimates}

Write \(C(A)=U_0 A U_0^*\), and \[
 B_j=C^j(B)=W(v_j)+W(-v_j),\qquad v_j=(qh^j,h^{-j}),
 \qquad \overline B=K^{-1}\sum_{j=0}^{K-1}B_j.                    \tag{FN.2}
\] The conjugation C shifts the label cyclically, so \(C(\bar B)=\bar B\). Define \[
 a_j=\left|e_N(q(h^j-h^{-j}))-1\right|,\qquad
 c_j=\frac{K-1-2j}{2K},\qquad G=\sum_{j=0}^{K-1}c_jB_j.          \tag{FN.3}
\] Then \(a_0=0,\ a_{K-j}=a_j\). Direct subtraction of the coefficients gives the exact cohomological equation \[
 G-C(G)=B-\overline B.                                         \tag{FN.4}
\] The coefficients have zero mean. Importantly, they are small at the half-period at which the large commutator phases accumulate.

Let \(P_{t}\) denote the second Pell sequence, \(\alpha^t=Q_t+P_t\sqrt2\). For the signed representative \(\delta\) of j mod K one has \[
 q(h^j-h^{-j})=-2P_{2\delta}\pmod N.                            \tag{FN.5}
\] Indeed \(h^\delta=(3+2s)^\delta=Q_{2\delta}+P_{2\delta}s\), and \(Fs=1\). For positive \(\delta=(K-1)/2-\ell\), \[
 \frac{P_{2\delta}}{N}
 =\frac{\lambda^{-\ell}}{\sqrt2\,\alpha}
       \frac{1-\lambda^{-2\delta}}{1-\lambda^{-K}}.              \tag{FN.6}
\] Put \(r=\lambda^{-1}\), \(A=4\pi/(\sqrt2\,\alpha)\). Using \(|e^{it}-1|\leq|t|\) and summing both signed representatives proves \[
 L_K:=\sum_{j=0}^{K-1}a_j
 \leq\frac{2A}{(1-r)(1-\lambda^{-K})}
 =\frac{2\sqrt2\pi}{1-\lambda^{-K}}<9\quad(K\geq3).              \tag{FN.7}
\] There is a sharper weighted estimate. At j=(K-1)/2-ell the weights at j and K-j have absolute values ell/K and (ell+1)/K. Therefore \[
 \sum_j |c_j|a_j
 \leq\frac{A}{K(1-\lambda^{-K})}\sum_{\ell\geq0}(2\ell+1)r^\ell
 =\frac{2\pi}{K(1-\lambda^{-K})}.                              \tag{FN.8}
\] The final identity follows from \(A(1+r)/(1-r)^2=2\pi\). The estimates are analytic inequalities; the regression's decimal phase values are only finite checks of them.

\subsection{Quadratic lifting and orthogonal commutator supports}\label{floquet:quadratic-lifting-and-orthogonal-commutator-supports}

\begin{lemma}[lifting]\label{pass1:ch07_floquet_resonance-2} Fix m. For all sufficiently large K, every signed m-term reciprocal relation \[
 \sum_{\nu=1}^m\epsilon_\nu h^{j_\nu}
 =\sum_{\nu=1}^m\epsilon_\nu h^{-j_\nu}=0\pmod N,
 \qquad \epsilon_\nu\in\{-1,1\},                               \tag{FN.9}
\] can be cyclically rotated to a polynomial \(P\) with \(P(\lambda)=0\). Its coefficient \(\ell^1\)-norm is at most \(m\), and its degree is at most \(K-\lceil K/m\rceil\). For \(m=4\), this holds for every odd prime \(K\geq11\), and the relation is then coefficientwise balanced.
\end{lemma}

\textbf{Proof.} Cancel opposite signs at identical labels. If anything remains, cut the cyclic support at a largest gap and rotate. There are at most \(m\) support points, giving the stated degree bound. Both congruences in (FN.9) survive. Divide \(P\) by \(Z^2-6Z+1\) over the integers. Its remainder \(aZ+b\) vanishes at \(h\) and \(h^{-1}\); their difference \(4s\) is a unit. Thus \(N\mid a\) and \(N\mid b\). A nonzero quadratic norm of \(P(\lambda)\) would be an integer divisible by \(N^2\), but \[
 |P(\lambda)P(\lambda^{-1})|
 \leq m^2\lambda^{K-\lceil K/m\rceil}<N^2                       \tag{FN.10}
\] for sufficiently large K at fixed m. This is impossible, so \(P(\lambda)=0\). For \(m=4\), a sufficient condition is \(\lambda^{K/4}(1-\lambda^{-K})^2>64\). It holds for \(K\geq11\): \(\lambda>5\), \(5^{11/4}>80\), and \(1-5^{-11}>9/10\) give a strict lower bound of 64.8. A nonzero polynomial of coefficient \(\ell^1\)-norm at most four cannot vanish at \(\lambda>3\), since its leading term strictly dominates the sum of all lower terms. Thus in this case the polynomial was zero. End of proof.

It follows that the \(2K(K-1)\) vectors \[
 \epsilon v_t+\delta v_s,\qquad 0\leq t<s<K,\quad
 \epsilon,\delta\in\{-1,1\},                                  \tag{FN.11}
\] are nonzero and distinct for \(K\geq11\). Any collision would give a balanced four-letter relation, forcing the original unordered signed pair. The matrices \(W(v)\) form an orthonormal basis for the full normalized Hilbert--Schmidt inner product, also for composite odd \(N\), because \(\operatorname{Tr}W(v)=N\mathbf1_{v=0}\).

Each of the four Weyl coefficients in \([B_t,B_s]\) has modulus \(a_{s-t}\). Different unordered time pairs have disjoint support by (FN.11). Since \(c_t-c_s=(s-t)/K\), exact orthogonality gives \[
 \|[G,\overline B]\|_{2,N}^2
 =\frac4{K^4}\sum_{j=1}^{K-1}(K-j)j^2a_j^2
 =\frac2{K^3}\sum_{j=1}^{K-1}j(K-j)a_j^2
 \leq\frac9K.                                                  \tag{FN.12}
\] Here \(\|A\|_{2,N}^2=\operatorname{Tr}(A^*A)/N\). The middle identity uses \(a_{K-j}=a_j\); the last uses \(j(K-j)\leq K^2/4\), \(a_j^2\leq2a_j\), and (FN.7).

The operator norm needs no orthogonality. Four Weyl coefficients, (FN.7), and (FN.8) give \[
 \|[G,B]\|_{\rm op}<\frac{26}{K},\qquad
 \|[B,\overline B]\|_{\rm op}<\frac{36}{K}.                     \tag{FN.13}
\] For a parity-preserving A, positivity of A*A implies \[
 \|A\|_{2,-}\leq\sqrt{\frac{2N}{N-1}}\,\|A\|_{2,N}.             \tag{FN.14}
\] No orthogonality of Weyl operators restricted to the odd subspace is assumed.

\subsection{A finite normal form without a small generator}\label{floquet:a-finite-normal-form-without-a-small-generator}

For finite Hermitian A,D, real \(\varepsilon\), and any unitarily invariant norm, \[
 \|e^{i\varepsilon A}e^{i\varepsilon D}
           -e^{i\varepsilon(A+D)}\|
 \leq\frac{\varepsilon^2}{2}\|[A,D]\|.                          \tag{FN.15}
\] This is the usual two-factor commutator estimate, not a new general product formula. Differentiate \(e^{i(\varepsilon-t)(A+D)}e^{itA}e^{itD}\). After factoring out unitaries, the integrand is \(e^{itA}De^{-itA}-D\). Integrate its derivative once more; the triangular integration domain has area \(\varepsilon^2/2\). Negative \(\varepsilon\) follows by replacing \(A,D\) with \(-A,-D\). In particular, \(|\varepsilon|\|A\|\) need not be small.

Set \(V_\varepsilon=e^{-i\varepsilon G}\). By (FN.4), \[
 V_\varepsilon U_\varepsilon V_\varepsilon^*
 =e^{-i\varepsilon G}e^{i\varepsilon B}
       e^{i\varepsilon C(G)}U_0.
\] Apply (FN.15) first to \(-G,B\), then to \(B-G,C(G)\). Their total sum is \(\bar B\), and \([B-G,C(G)]=[B-G,\bar B]\). Equations (FN.12)--(FN.14) give \[
 \begin{split}
 \|V_\varepsilon U_\varepsilon V_\varepsilon^*
       -e^{i\varepsilon\overline B}U_0\|_{2,-}
 &\leq\frac{\varepsilon^2}{2}
   \left(\frac{62}{K}+3\sqrt{\frac{2N}{(N-1)K}}\right)\\
 &\leq\frac{12\varepsilon^2}{\sqrt K},\qquad K\geq11.
 \end{split}\tag{FN.16}
\] For the last step it even suffices to use \(N\geq7\) and \(31/\sqrt{11}+(3/2)\sqrt{7/3}<12\). This proves (FN.D) for all real \(\varepsilon\). G may have operator norm of order K. Expanding its exponential in powers of \(\varepsilon\) and requiring \(\varepsilon K\) small would incorrectly discard the critical regime. Estimate (FN.15) avoids that requirement.

Since \(\bar B\) commutes with \(U_{0}\) and \(U_0^K=I\), telescoping unitary powers in the same norm yields, for any integer m, \[
 \left|\tau_-(U_\varepsilon^{mK})
       -\tau_-(e^{im\varepsilon K\overline B})\right|
 \leq12|m|\varepsilon^2\sqrt K,\qquad
 \tau_-(A)=d_N^{-1}\operatorname{Tr}_{-}A.                      \tag{FN.17}
\] For negative m use adjoints. Cauchy--Schwarz for the normalized trace justifies the passage from the norm to traces.

\subsection{A self-contained global law for the averaged observable}\label{floquet:a-self-contained-global-law-for-the-averaged-observable}

The remaining input is established here directly on the whole odd subspace; no unproved claim about the old arithmetic-moment corpus is imported. Define \[
 X_K=\sqrt{K/2}\,\overline B=(2K)^{-1/2}\sum_{j=0}^{K-1}B_j.
\]

\begin{lemma}[Gaussian average]\label{pass1:ch07_floquet_resonance-3} For each fixed integer \(r\geq1\), \[
 \tau_-(X_K^{2r})=(2r-1)!!+O_r(K^{-1})+O_r(K^r/N),               \tag{FN.18}
\] and for each fixed odd m, eventually \[
 |\tau_-(X_K^m)|\leq (2K)^{m/2}/(N-1).                         \tag{FN.19}
\] The empirical measures of \(X_{K}\) converge weakly and in quadratic Wasserstein distance to the standard normal law.
\end{lemma}

\textbf{Proof.} In the full Weyl representation \(\operatorname{Tr}(RW(v))=1\) for every v: the unique fixed point \(x=-b/2\) of RW(a,b) contributes phase 1. Thus the normalized odd trace of a word \(e_{N}\)(\(\phi\))W(v) is \[
 e_N(\phi)\frac{N\mathbf1_{v=0}-1}{N-1}.                       \tag{FN.20}
\] Expanding \(X_K^m\), the noncentral part is bounded by \((2K)^{m/2}/(N-1)\). The central part is the weighted count of the reciprocal relations (FN.9), normalized by \([N/(N-1)](2K)^{-m/2}\).

Here is the count including structurally forced nonbalanced relations. For a polynomial of coefficient \(\ell^1\)-norm at most m vanishing at \(\lambda\), a gap \(d>2\log_\lambda m\) separates two independently vanishing pieces. Write \(P=P_{\rm low}+Z^{a+d}P_{\rm high}\), with \(\deg P_{\rm low}\leq a\) and nonnegative exponents in \(P_{\rm high}\). If \(P_{\rm high}(\lambda)\ne0\), its algebraic-integer norm is a nonzero integer. Since \(|P_{\rm high}(\lambda^{-1})|\leq m\), we have \(|P_{\rm high}(\lambda)|\geq1/m\). Cancellation with \(|P_{\rm low}(\lambda)|\leq m\lambda^a\) would force \(\lambda^d\leq m^2\), a contradiction.

Repeatedly splitting gives components with span bounded in terms of m, for example by \(D_m=(m-1)(\lceil2\log_\lambda m\rceil+1)\). Every nonzero component has at least three signed occurrences: two distinct positive powers cannot cancel with signs of magnitude one. Restore previously cancelled opposite pairs as two-occurrence components. A nonbalanced 2r-letter word therefore has at most r-1 components. For each of finitely many partitions of the 2r positions, signs and bounded relative offsets, each component has at most K translations on the cyclic time set. There are only \(O_r(K^{r-1})\) such words. There is no additional factor K from choosing a cut: the count is over the component translations themselves, and the largest-gap cut was used only to prove bounded width. Their normalized contribution is \(O_r(1/K)\).

For odd m, a lifted P is an integer multiple of \(Z^2-6Z+1\), so \(P(1)\) is divisible by -4. But \(P(1)\) is the sum of an odd number of signs. Thus no central word exists once Lemma~\ref{pass1:ch07_floquet_resonance-2} applies, proving (FN.19).

A balanced \(2r\)-word pairs identical time labels with opposite signs. Assignments with repeated pair labels cause only \(O_r(K^{r-1})\) overcount. Each matching of the \(2r\) positions has \(K^r\) pair-label choices and \(2^r\) orientations. Its Weyl phase is a product, over crossing pairs, of factors \(e_N(q(h^a-h^{-a}))\) or their conjugates. Summing the signed symplectic products in word order shows that disjoint or nested pairs cancel, while a crossing pair retains one full symplectic product.

Each crossing difference of independent uniform pair labels is uniform modulo \(K\). Independence of the different crossing factors is neither true in general nor needed. The inequality \(|\prod_l u_l-1|\leq\sum_l|u_l-1|\), for \(|u_l|=1\), and (FN.7) give average phase error at most nine times the number of crossings divided by \(K\). Hence each matching contributes \(1+O_r(1/K)\) after normalization. There are \((2r-1)!!\) matchings. Together with the nonbalanced count and (FN.20), this proves (FN.18).

The second-moment bound gives tightness. Higher even moments give uniform integrability when passing moments to subsequential limits. The Gaussian moment sequence is determinate, for example by Carleman\textquotesingle s criterion, since \(((2r-1)!!)^{1/(2r)}\leq\sqrt{2r}\). Every subsequential limit is standard normal. Convergence of the second moments upgrades weak convergence to quadratic Wasserstein convergence. End of proof.

\subsection{Critical and subcritical propagator spectra}\label{floquet:critical-and-subcritical-propagator-spectra}

For \(\varepsilon=\tau/\sqrt{2K}\), (FN.17) becomes \[
 \left|\tau_-(U_\varepsilon^{mK})-
              \tau_-(e^{im\tau X_K})\right|
       \leq \frac{6|m|\tau^2}{\sqrt K}.                        \tag{FN.21}
\] Lemma~\ref{pass1:ch07_floquet_resonance-3} and bounded continuity of exp(im \(\tau\) x) give (FN.A). Integer Fourier coefficients determine measures on the circle: trigonometric polynomials are dense in continuous functions there. Applying the characteristic function of a normal variable proves (FN.B), and taking absolute squares proves (FN.C). No eigenphase labelling across the old degenerate sectors is used or asserted.

There is also a genuine extension of the earlier subcritical global transfer, but for \textbf{folded phases} rather than individually isolated clusters. Let \(\varepsilon_{K}\) be nonzero and \(\tau_K=\varepsilon_K\sqrt{2K}\) tend to zero. If \(e^{i\theta_l}\) are the eigenvalues of \(U_{\varepsilon,K}\), define \[
 y_l=\frac{\operatorname{Arg}(e^{iK\theta_l})}{\tau_K},
 \qquad \operatorname{Arg}\in[-\pi,\pi).
\]

\begin{corollary}[subcritical folded Gaussian law]\label{pass1:ch07_floquet_resonance-4} The empirical measures of \(y_{l}\) converge in quadratic Wasserstein distance to the standard normal law whenever \(\varepsilon_K\sqrt K\) tends to zero.
\end{corollary}

\textbf{Proof with a quantitative tail control.} For two unitary matrices, the squared absolute overlaps of their orthonormal eigenbases form a doubly stochastic matrix. The corresponding average squared chord distance equals their normalized squared Hilbert--Schmidt difference. Decomposing this matrix into permutation matrices gives an eigenvalue matching with no larger average cost. The decomposition can be obtained by Hall's matching criterion and successive subtraction; it uses no eigenvalue simplicity. Thus (FN.16), telescoped K times, matches the powered spectrum to \(e^{i\tau_Kx_l}\), where \(x_{l}\) are the eigenvalues of \(X_{K}\), with root mean square chord error at most \(6\tau_K^2/\sqrt K\).

If \(|\tau_Kx_l|\leq\pi/2\), principal argument differences obey \(|\operatorname{Arg}u-\tau_Kx_l|\leq4|u-e^{i\tau_Kx_l}|\). On the complementary set, \(|\operatorname{Arg}u/\tau_K-x_l|\leq3|x_l|\). Consequently \[
 \begin{split}
 W_2(\operatorname{law}(y),\operatorname{law}(X_K))
 &\leq\frac{24|\tau_K|}{\sqrt K}
   +3\sqrt{\tau_-\!\left(X_K^2
       \mathbf1_{|X_K|>\pi/(2|\tau_K|)}\right)}\\
 &\leq\frac{24|\tau_K|}{\sqrt K}
       +\frac{6|\tau_K|}{\pi}\sqrt{\tau_-(X_K^4)}.
 \end{split}\tag{FN.22}
\] The fourth moment stays bounded by Lemma~\ref{pass1:ch07_floquet_resonance-3}, so the right side tends to zero. Combine with the \(W_{2}\) conclusion of that lemma. End of proof.

The former isolated-cluster argument required \(\varepsilon_KK^{3/2}\) tending to zero. Corollary~\ref{pass1:ch07_floquet_resonance-4} permits an additional factor K in the kick scale, but its observable is a folded spectral measure. It must not be restated as an individual-cluster rank or no-overlap theorem.

\subsection{A compact-test microscopic obstruction}\label{floquet:a-compact-test-microscopic-obstruction}

A negative statement about microscopic correlations needs less information than a positive universality theorem. The recurrence law above supplies such a statement through Fourier positivity. This section concerns the whole odd spectrum, with its multiplicities and the \textbf{global mean-density unfolding}. It is not a claim about a further symmetry sector or an unfolding which first removes a nonconstant mesoscopic intensity profile.

Let \(x_{1}\),\ldots,\(x_{d}\) be the eigenphases of a d-dimensional unitary U, measured in turns, so its eigenvalues are \(e^{2\pi ix_j}\). For a compactly supported continuous test \(\phi\) define the ordered factorial pair statistic \[
 C_d(U;\phi)=\frac1d\sum_{a\ne b}\sum_{n\in\mathbb Z}
              \phi\bigl(d(x_a-x_b+n)\bigr).                    \tag{FN.23}
\] Distinct indices, not distinct eigenvalues, are required in this definition. For \(L>0\) use the compact triangular test \[
 \phi_L(s)=\frac1L(1-|s|/L)_+,\quad
 \int\phi_L=1,\quad \phi_L(0)=1/L,\quad
 \widehat\phi_L(u)=\operatorname{sinc}^2(Lu)\geq0,               \tag{FN.24}
\] where \(\operatorname{sinc}(v)=\sin(\pi v)/(\pi v)\), with value 1 at zero. Fourier transforms use \(e^{-2\pi ius}\). Absolute convergence of the periodized Fourier series (the coefficients decay quadratically) gives the exact identity \[
 C_d(U;\phi_L)=d^{-2}\sum_{t\in\mathbb Z}
   \widehat\phi_L(t/d)(|\operatorname{Tr}U^t|^2-d).              \tag{FN.25}
\] If \(d>L\), the total diagonal subtraction is exactly 1/L, because \(\sum_n\phi_L(dn)=\phi_L(0)\). For any distinct positive integers \(r_{1}\),\ldots,\(r_{M}\), positivity therefore gives \[
 C_d(U;\phi_L)\geq1-\frac1L+
   2\sum_{j=1}^M\operatorname{sinc}^2(Lr_j/d)
                   |d^{-1}\operatorname{Tr}U^{r_j}|^2.         \tag{FN.26}
\] No control or truncation of the remaining high-frequency traces is needed for this lower bound; their contribution is nonnegative.

There is an equivalent counting interpretation. Let \(N_{L}\)(y) count the phases in the circular interval {[}y,y+L/d), with y uniform in {[}0,1). Then \(E_{y}\) \(N_{L}\)=L and an overlap-of-intervals calculation gives \[
 \operatorname{Var}_y N_L
      =L^2\bigl(C_d(U;\phi_L)-1\bigr)+L.                       \tag{FN.27}
\] This identity is also verified independently by integrating the piecewise-constant counting function, rather than reusing the pair sum.

\begin{theorem}[obstruction and necessary kick scale]\label{pass1:ch07_floquet_resonance-5} For the kicked Pell family of Theorem~\ref{pass1:ch07_floquet_resonance-1}, put \(\tau_K=\varepsilon_K\sqrt{2K}\). If \(\tau_{K}\) tends to a finite nonzero \(\tau\), then for every fixed \(L>0\), \[
 \liminf_{K\to\infty} C_{d_N}(U_{\varepsilon_K,N};\phi_L)
      \geq \Theta(\tau)-\frac1L,\qquad
 \Theta(\tau)=\sum_{m\in\mathbb Z}e^{-m^2\tau^2}>1,             \tag{FN.28}
\] and \[
 \liminf_{K\to\infty}\operatorname{Var}_y N_L
      \geq L^2\bigl(\Theta(\tau)-1\bigr).                      \tag{FN.29}
\] If \(\tau_K\to0\), then \(C_{d_N}(U_{\varepsilon_K,N};\phi_L)\to+\infty\) for every fixed \(L\). Consequently, convergence of the globally unfolded pair measures to a bounded density \(\rho(s)\) with \(\rho(s)\to1\) as \(|s|\to\infty\) requires \[
             |\varepsilon_K|\sqrt K\longrightarrow\infty.     \tag{FN.30}
\] This is a necessary condition only, not a sufficient universality criterion.
\end{theorem}

\textbf{Proof.} The recurrence trace limit remains valid when \(\tau_K\to\tau\) varies with \(K\). Lemma~\ref{pass1:ch07_floquet_resonance-3} bounds \(\tau_-(|X_K|)\), and \(|\tau_-(e^{im\tau_KX_K})-\tau_-(e^{im\tau X_K})|\leq |m|\,|\tau_K-\tau|\,\tau_-(|X_K|)\). The error in (FN.21) is uniform for bounded \(\tau_K\). Thus the normalized \(mK\)-trace tends to \(e^{-m^2\tau^2/2}\) for each fixed \(m\), including \(\tau=0\).

Apply (FN.26) with \(r_j=jK\). Since \(K/d_N\to0\), the sinc factors tend to one. At each fixed \(M\), the resulting lower bound is \(1-1/L+2\sum_{m=1}^M e^{-m^2\tau^2}\). Let \(M\) tend to infinity \emph{after} taking the lower limit in \(K\). This proves (FN.28), or divergence when \(\tau=0\). Equation (FN.27) gives (FN.29).

Suppose the pair measures converge vaguely to \(\rho\)(s)ds, where \(\rho\) is bounded and clusters to unit intensity. For every fixed L the limiting triangular statistic would be \[
 A_L=\int_{-1}^1(1-|u|)\rho(Lu)\,du.
\] Dominated convergence gives \(A_L\to1\) as \(L\to\infty\). But (FN.28) requires \(A_L\geq\Theta(\tau)-1/L\) for every \(L\), a contradiction. For \(\tau=0\), a finite locally bounded target already contradicts the fixed-\(L\) divergence. If \(|\varepsilon_K|\sqrt K\) does not tend to infinity, a bounded subsequence of \(\tau_K\) has a further subsequence converging to a finite limit. The same contradiction applies. End of proof.

The use of a compact test is essential for the logical scope. A Gaussian test without an independently controlled tail would only obstruct a Schwartz-test limit, not automatically vague local convergence. If the pair convergence definition uses only smooth compact tests, positive measures acquire locally uniform mass bounds by testing a bump which is one on a given compact set; approximation then extends convergence to the continuous triangular tests. Alternatively, start directly with a smooth nonnegative compact f of integral one, take \(\phi\)=f*tilde(f), and scale it: \(\widehat\phi=|\widehat f|^2\geq0\) and the same proof uses \(\phi\)(0)/L instead of 1/L. No unproved high-frequency-tail assumption is hidden in Theorem~\ref{pass1:ch07_floquet_resonance-5}.

The Poisson pair density \(\rho\)=1 and the sine-kernel density \(\rho_2(s)=1-\operatorname{sinc}^2(s)\) are both excluded. More explicitly, keeping only the first resonance gives \[
 \liminf C_{d_N}(U_{\varepsilon_K,N};\phi_L)
        \geq1+2e^{-\tau^2}-1/L.                               \tag{FN.31}
\] Since \(\int\phi_L\rho_2\leq1\), taking \(L=e^{\tau^2}\) yields a gap of at least \(e^{-\tau^2}\) from the sine-kernel prediction. At \(\tau=1\), already L=2 gives lower bound 1.2357588823\ldots, greater than both that prediction and the Poisson value 1.

The standard GOE bulk pair density is excluded by the bounded-clustering part of the theorem, not by an assumed pointwise bound \(\rho\leq1\). In unit intensity normalization it is the even continuation of \[
 \rho_1(s)=1-S(s)^2-S'(s)\int_s^\infty S(t)\,dt,\qquad
 S(s)=\operatorname{sinc}(s),\quad s\geq0.                       \tag{FN.32}
\] This standard formula is recorded in Brezin--Hikami, equation (181) after unfolding. It is bounded at zero, and integration by parts in the sine tail shows \(\rho_1(s)=1+O(s^{-2})\) at infinity. The same elementary criterion therefore applies. The random-matrix formulas, Fourier positivity and counting-variance identity are standard ingredients; no priority is claimed for those general tools.

This changes the route toward a positive BGS result: the critical and subcritical globally unfolded odd spectra considered here are not merely unproved universal families; they are \textbf{excluded} from the bounded clustering pair limits just specified. The statement does not settle nearest-neighbour distributions individually, density-adapted re-unfolding, a further desymmetrized block, or fixed nonzero classical kicks. It is not a refutation of a correctly qualified generic BGS conjecture.

The finite verifier keeps forty configurations (K=3,5,7, several \(\tau\) and L) and two synthetic controls. At K=7,\(\tau=1\),L=2 it numerically gives \(C_{d}\)=1.30936229449\ldots{} and Var \(N_{L}\)=3.23744917797\ldots, while the finite lower bound retaining K and 2K is 1.21861711382\ldots{} . These are floating-point checks of finite identities, not the proof of (FN.28). In particular the asymptotic bound is not incorrectly enforced on small K.
